% ============================================================================
%  it_from_bit.tex
%
%  "It from Bit: A Concrete Attempt"
%  A toy universe built from a bit-local lattice automaton.
%
%  Compile with:  pdflatex it_from_bit.tex   (run twice for references)
% ============================================================================

\documentclass[11pt,a4paper]{article}

\usepackage[T1]{fontenc}
\usepackage[utf8]{inputenc}
\usepackage[english,brazilian]{babel}
\usepackage{amsmath,amssymb}
\usepackage[margin=2.6cm]{geometry}
\usepackage[hidelinks]{hyperref}

\newcommand{\code}[1]{\texttt{#1}}

\title{\textbf{It from Bit: A Concrete Attempt}\\[2mm]
       \large Emergent physics from a bit-local lattice automaton}

\author{The Alpha Project\thanks{Source code of the \emph{Toy Universe}
        automaton: the \texttt{alpha} repository.}}
\date{\today}

\begin{document}

\maketitle

\begin{abstract}
John Archibald Wheeler's dictum ``it from bit'' holds that every physical
object---every \emph{it}---derives its existence, its name, and its laws from
bits, from information-theoretic answers to yes-or-no questions
\cite{wheeler1990,wheeler1989}. This manuscript presents a concrete, fully
executable attempt at that program: a \emph{toy universe} in which geometry,
a wave-like light clock, charge, polarization, and a matter--antimatter
asymmetry all emerge as degrees of freedom of a single deterministic lattice
automaton. The model is a 3-torus cellular automaton whose cells carry six
charge bits updated by integer arithmetic only, organized into winding layers
with an island-affinity structure. A pulsating spherical wavefront---a
constant-speed ``light clock''---schedules all local dynamics; an
election--broadcast--reconstruction protocol lets a transverse polarization
pair $(u,v)$ emerge from pure bit propagation; and a capture/escape balance at
the level of islands realizes a dynamic charge-quantization attractor. The
automaton is instrumented end to end: shell completeness, radial amplitude
profiles, charge census, and island-population statistics are sampled every
light frame. We argue that the construction is a faithful, if miniature,
instance of the \emph{it from bit} thesis, and we discuss which features are
model-dependent engineering and which hint at genuinely information-theoretic
origins.
\end{abstract}

\begin{otherlanguage}{brazilian}
\begin{abstract}
O ditado de John Archibald Wheeler, \emph{it from bit} (``o isto vem do
bit''), sustenta que todo objeto físico---todo ``isto''---deriva sua
existência, seu nome e suas leis de bits, das respostas informacionais a
perguntas de sim ou não. Este manuscrito apresenta uma tentativa concreta e
plenamente executável desse programa: um universo de brinquedo no qual a
geometria, um relógio de luz ondulatório, a carga, a polarização e uma
assimetria matéria--antimatéria emergem como graus de liberdade de um único
autômato celular determinístico. O modelo é um autômato celular sobre um
3-toro cujas células carregam seis bits de carga, atualizadas apenas com
aritmética inteira, organizadas em camadas de enrolamento com uma estrutura
de afinidade por ilhas. Uma frente de onda esférica pulsante---um ``relógio
de luz'' de velocidade constante---agenda toda a dinâmica local; um protocolo
de eleição--difusão--reconstrução faz emergir um par de polarização
transversal $(u,v)$ a partir da pura propagação de bits; e um balanço
captura/escape ao nível das ilhas realiza um atrator de quantização dinâmica
de carga. O autômato é instrumentado de ponta a ponta: completude de casca,
perfis radiais de amplitude, censo de carga e estatísticas de população de
ilhas são amostrados a cada quadro de luz. Argumentamos que a construção é
uma instância fiel, ainda que miniatura, da tese do \emph{it from bit}, e
discutimos quais características são engenharia dependente do modelo e quais
sugerem origens genuinamente informacionais.
\end{abstract}
\end{otherlanguage}

\noindent\textbf{Keywords:} it from bit, digital physics, cellular automata,
emergent spacetime, emergent charge, wavefront automaton, computational
universe.

\section{Introduction}
\label{sec:intro}

``It from bit---otherwise put, every \emph{it}---every particle, every field
of force, even the spacetime continuum itself---derives its function, its
meaning, its very existence \ldots{} from the answers to yes-or-no questions
\ldots{} from bits'' \cite{wheeler1990,wheeler1989}. With these words John
Archibald Wheeler posed the strongest form of the thesis that physics is, at
bottom, information processing. The slogan has since become the banner of
\emph{digital physics}: the conjecture that the universe is a discrete,
programmable system---a computer of sorts---and that its apparent continuity
is an emergent fiction \cite{zuse1967,fredkin1990,wolfram2002,
toothooft2016,zenil2012}.

The program is easy to state and notoriously hard to carry out. Most
proposals remain at the level of metaphor: a slogan, a research programme, a
dictionary between information-theoretic and physical vocabulary, but no
executable object whose dynamics actually \emph{produce} particles, charges,
or geometry. This manuscript is a deliberate, engineering-flavored
counterpoint. We present a complete, running, fully instrumented
automaton---a \emph{toy universe}---in which the central items of the physics
vocabulary are \emph{read off} a single bit-local update rule rather than
postulated as fields.

The construction follows five design choices, each intended to be as concrete
as possible.

\begin{enumerate}
  \item \textbf{Everything is a bit.} The fundamental state is a lattice of
        cells, each carrying six charge bits and a handful of auxiliary
        flags. There are no real-valued fields: every update is an integer
        operation.
  \item \textbf{Geometry is computed, not stored.} Squared distances and
        integer radii from each source center are maintained by an
        incremental, add-only integer algorithm, so a spherical shell is an
        emergent set of cells rather than a stored mathematical surface.
  \item \textbf{Time is a local light clock.} A pulsating wavefront---an
        expanding and contracting spherical shell of period $2R_{\max}$---
        replaces the continuous time parameter of field theory with a
        discrete, locally available signal.
  \item \textbf{Structure emerges through broadcast.} The transverse
        polarization pair $(u,v)$ and the momentum axis are not stored state:
        every cell \emph{reconstructs} them from a single scalar ``arrival
        stamp'' broadcast by a helical walker.
  \item \textbf{Everything is measured.} The automaton samples its own
        observables---shell completeness, radial amplitude profiles, charge
        census, island populations---every light frame, yielding quantitative
        output comparable against the expectations of a constant-speed wave.
\end{enumerate}

The remainder of this paper is organized as follows. Section~\ref{sec:model}
defines the automaton. Section~\ref{sec:instrumentation} describes the
built-in instrumentation. Section~\ref{sec:implementation} summarizes the
implementation. Section~\ref{sec:results} reports representative
observations. Section~\ref{sec:discussion} discusses, in Wheeler's spirit,
what the toy universe can and cannot teach us about ``it from bit'' as a
physical thesis. Section~\ref{sec:conclusion} concludes.

\section{The model}
\label{sec:model}

\subsection{Lattice, cell state, and update schedule}
\label{sec:lattice}

The arena is a three-torus of edge length $L$: coordinates
$x,y,z\in\{0,\dots,L-1\}$ wrap around, so a step across a face of the cube
continues from the opposite face as a pure translation---orientation is
preserved, nothing is mirrored. The automaton in fact allocates a stack of
$W$ such lattices, indexed by a fourth coordinate $w$, the \emph{winding}
index. Cross-layer adjacency is governed not by spatial geometry but by a
rotation schedule (a rotate-partners operation), so the $w$ axis is a discrete
analogue of an internal (winding) degree of freedom.

Each cell of the combined lattice of size $L^3 \times W$ is a compact record:

\begin{itemize}
  \item six \emph{charge bits}---three colour bits $c_0,c_1,c_2$, one bit $q$,
        and two weak bits $w_0,w_1$---packed into a single byte, with masks
        $C_0=0\times01$, $C_1=0\times02$, $C_2=0\times04$, $Q=0\times08$,
        $W_0=0\times10$, $W_1=0\times20$;
  \item an intrinsic identity $w$ and an auxiliary \emph{leader} identity
        shared by an entire winding island;
  \item an \emph{affinity} label $a$ fixing island membership, with a special
        value marking orphan cells;
  \item integer squared radius $r^2$ and integer radius $r$ from the layer's
        moving source center;
  \item the phase flags $p_B$ and $s_B$ (local wave sign and emergent
        transverse flag);
  \item the emergent polarization pair $(u,v)$ and the broadcast arrival
        stamp $b$.
\end{itemize}

Every tick each cell is updated from its neighbourhood by the composition of
elementary operations---convolution (local averaging of the phase signal),
diffusion (spreading of affinity), and relocation (reissue of released
charges)---optionally with programmable delays between the phases. All
operations are bit and integer operations; floating point never appears in
the update kernel.

\subsection{The pulsating wavefront: a lattice light clock}
\label{sec:wavefront}

A single master schedule drives the entire automaton. Each winding layer owns
a source center; the active shell radius is the triangle wave

\begin{equation}
  \label{eq:lightclock}
  \rho(t) \;=\; \begin{cases}
    \tau, & \tau \le R_{\max},\\
    2R_{\max} - \tau, & \tau > R_{\max},
  \end{cases}
  \qquad \tau = t \bmod 2R_{\max},
\end{equation}

with period $2R_{\max}$ and amplitude $R_{\max}=L/2$. A cell is
\emph{active} exactly when its propagated integer radius $r$ equals
$\rho(t)$; because the peak advances one cell per tick, the front is a
constant-speed wave---a lattice-native ``light clock'' that, unlike a global
integer time counter, is available locally to every cell. At the turnaround
$t \equiv R_{\max}$ the shell stops, inverts, and \emph{elects} a direction:
this tick is the dynamical event on which the emergent structure of
Section~\ref{sec:polarization} hangs.

The radius field itself is never recomputed by square roots. It is maintained
by an incremental CaRaSh-style distance transform: after each step of the
source center, every cell updates $r^2$ and $r$ using additions and simple
square-boundary tests, so the shell ``just happens'' as the collection of
cells whose radius matches the schedule~\eqref{eq:lightclock}.

\subsection{Winding layers and island affinity}
\label{sec:islands}

Winding layers are grouped into \emph{islands} of fixed size $3L^2$; the
chief layer of an island is its first member. Initial conditions place, in
each layer, a sphere of radius $R_{\max}$ around the layer center; cells
inside the sphere inherit the island chief as affinity, cells outside are
orphans. Affinity is dynamic: across light frames, \emph{capture} events turn
orphans into members, \emph{escape} events release members back into the
orphan pool, and a member changing island counts as escape plus capture.
These events are the fuel of the charge-quantization attractor of
Section~\ref{sec:attractor}.

\subsection{Emergent polarization pair \texorpdfstring{$(u,v)$}{(u,v)}}
\label{sec:polarization}

Perhaps the most distinctive mechanism is the way the model refuses to store
a polarization field. Instead, a three-stage protocol reconstructs it from a
scalar bit signal.

\textbf{Election.} At every expansion limit of the breathing wavefront---the
exact tick at which the wave turns from ascending to descending, i.e.\
$t\equiv R_{\max}$---every active shell cell is seeded with the payload
\begin{equation}
  \label{eq:election}
  \mathrm{payload}(x) = \big(\mathrm{score}(x) \ll 24\big) \mid
  \mathrm{code}(x),
  \qquad
  \mathrm{score}(x) = H\big((\nu(x)+s)\bmod 9L,\;\mathrm{code}(x)\big),
\end{equation}
where $\nu(x)$ is the island index of the cell, $\mathrm{code}(x)$ its linear
index, $H$ a fixed integer hash and $s$ a global seed dial. Because the code
occupies the low bits, the payload order is a strict total order: ties are
impossible, so the winner is unique and is identified once, at sowing time.
The displacement from the lattice center to the winning cell becomes the
momentum vector $\mathbf{m}$, rescaled to $|\mathbf{m}|=L/2$.

\textbf{Broadcast.} A helical walker traces a spiral around the cylinder whose
axis is the elected direction, advancing one cell per tick using only
additions, subtractions, shifts, and comparisons among its six face
neighbours. Each freshly computed spiral point stamps the current tick into
the cell's arrival field $b$ of a second value lattice (zero meaning ``never
reached''). Every other tick, with alternating sweep directions, each cell
copies the greatest neighbour stamp whenever it exceeds its own, so the news
diffuses to every cell in finite time.

\textbf{Reconstruction.} The arrival stamp is a phase source:
\begin{equation}
  \label{eq:reconstruction}
  \phi_{\mathrm{full}} = 2R^2,\qquad R = R_{\max}-2,
  \qquad
  \phi = (b-1)\bmod \phi_{\mathrm{full}},
  \qquad
  j = \lfloor \phi / R \rfloor,
\end{equation}
and the transverse pair is
\begin{equation}
  \label{eq:u}
  u = \begin{cases} R(R-2j), & j < R,\\ R(2j-3R), & j \ge R, \end{cases}
  \qquad
  v = \begin{cases} 2R\,\mathrm{isqrt}\!\big(j(R-j)\big), & j < R,\\[2mm]
       -2R\,\mathrm{isqrt}\!\big((j-R)(2R-j)\big), & j \ge R, \end{cases}
\end{equation}
where $\mathrm{isqrt}$ is the integer square root. The pair
\eqref{eq:u} approximates the circle $u^2+v^2=R^4$, exactly when the integer
square root is exact. Polarization is thus not a field that evolves; it is a
bit that \emph{travels}, and a rule that \emph{decodes}.

\subsection{Charge dynamics and matter--antimatter census}
\label{sec:charges}

The charge byte is the only particle-like memory of a cell. Its six bits are
recombined by the local update rule as islands interact; the rule is written
so that the $q$ bit and the weak bits participate in the capture/escape
balance that re-emits charge at each island turnaround (``clock reset''). A
separate, read-only instrumentation layer maintains a per-tick census of
matter and antimatter populations, split by the two global sectors
(\emph{Orbis} and \emph{Umbra}), together with a ledger of ``virgin
wraps''---cells whose winding parity has never flipped. The census is purely
observational: by construction it never writes to the lattice, so the
automaton's dynamics are bit-identical with or without it.

\section{Instrumentation and observables}
\label{sec:instrumentation}

A toy universe is only scientific if it can be measured. The automaton ships
with four instrument modules, all reading \emph{only} the committed state of
the lattice and all emitting machine-readable reports.

\paragraph{Wavefront fidelity.}
Once per completed light frame, for each integer radius $r$ that was the
active target of some layer, the observed number of active cells is divided
by the theoretical number of lattice points with
$r^2 \le \Delta x^2+\Delta y^2+\Delta z^2 < (r+1)^2$ (exact integer counting,
no square roots). A faithful constant-speed front should reach completeness
$\approx 1$. The module also records the radial amplitude profile
$\langle |u| \rangle(r)$ correlated against $\sin(r)/r$, and the peak-radius
error---the distance between the observed peak of $\langle|u|\rangle$ and the
target radius---a direct check that the wave peak advances one cell per tick.

\paragraph{Charge census.}
The matter/antimatter census of Section~\ref{sec:charges}, throttled per
tick, with capture and escape rates per island.

\paragraph{Dynamic charge-quantization attractor.}
\label{sec:attractor}
For every island, the population $N(t)$ is sampled each light frame, together
with gross capture and escape counts. The hypothesis, stated in the model as
a target to be tested, is that there is a stable attractor $N^{*}$:
\begin{equation}
  \label{eq:attractor}
  \Gamma_{\mathrm{cap}}(N) > \Gamma_{\mathrm{esc}}(N)\ \text{for small }N,
  \qquad
  \Gamma_{\mathrm{cap}}(N) < \Gamma_{\mathrm{esc}}(N)\ \text{for large }N,
\end{equation}
with $\Gamma_{\mathrm{cap}}(N^{*}) \approx \Gamma_{\mathrm{esc}}(N^{*})$. The
module performs an OLS fit of $\Delta N(t)=N(t+1)-N(t)$ against $N(t)$; a
negative slope is a restoring coefficient, and $N^{*}$ is estimated as
$-\mathrm{intercept}/\mathrm{slope}$, with standard errors. Time series can
be saved to and restored from disk, and written as CSV for external analysis.

\paragraph{Tomography.}
A visualization instrument slices, inverts, and animates the volume along any
axis with adjustable thickness, so that interior structure---shells, elected
axes, island shapes---can be inspected directly.

\section{Implementation}
\label{sec:implementation}

The automaton is written in C++17. The simulation runs on a dedicated thread;
rendering and interaction run on the main thread, and the voxel buffer is
guarded by a mutex so the two never race. A condition variable synchronizes
thread readiness; all cross-thread flags are atomic. The renderer uses
OpenGL~4.6 core profile through GLFW, GLAD and GLM, with a full HUD, menus,
recorder/replay pipeline, and a 2D/3D projection manager. An optional CUDA
bridge mirrors the three lattices on the device and accelerates the update
kernel when a GPU is available; a host-side stub preserves identical semantics
otherwise. All run parameters---lattice size, scenario, delays, field
visualization toggles, projection, tomography settings---are read from a
plain-text configuration file at startup, so every reported number is
reproducible.

\section{Results}
\label{sec:results}

The instrumentation is designed so that three claims can be checked
quantitatively at runtime.

\textbf{The front is a faithful light clock.} Shell completeness near unity
and a peak-radius error of at most one cell per light frame confirm that the
integer distance transform and the triangle schedule
\eqref{eq:lightclock} really implement a constant-speed front, and that the
``speed of light'' is a property of the rule, not of a stored time parameter.

\textbf{Polarization is a travelling bit.} The reconstruction
\eqref{eq:reconstruction}--\eqref{eq:u} predicts that every cell that ever
receives a nonzero arrival stamp lies, up to integer-square-root error, on
the circle $u^2+v^2=R^4$; monitoring the distribution of $u^2+v^2-R^4$ over
the lattice converts ``emergent polarization'' into a testable bound on the
automaton.

\textbf{Charge quantizes.} A negative fitted slope in
\eqref{eq:attractor} with a well-constrained $N^{*}$ demonstrates that the
capture/escape balance self-organizes a preferred charge population per
island---a genuinely emergent conservation law, in the sense that no bit of
the rule states it explicitly.

A full quantitative report, with per-run tables and figure panels produced
by the headless runner, is the subject of a companion document; this draft
concentrates on the architecture and on the falsifiable statements the
instrumentation exists to test. We stress what the toy universe does
\emph{not} claim: it does not quantize a pre-existing field, it does not
store a metric, and it does not invoke randomness---it is deterministic and
scale-free. Its virtue is that every emergent word in the previous paragraphs
has a precise, executable meaning.

\section{Discussion}
\label{sec:discussion}

Wheeler asked whether physics might be, at bottom, a matter of ``bits
answering yes-or-no questions.'' The toy universe is a concrete laboratory
for that question, and its architecture exposes three lessons.

First, \emph{emergence without a substrate}. Charge, momentum, and
polarization here are not fields that live on a pre-existing geometry; they
are decoded from propagation histories on a bit lattice whose geometry is
itself maintained by an integer transform. There is no level at which one is
forced to introduce a real number.

Second, \emph{the price of emergence is instrumentation}. A polarization
pair that is reconstructed rather than stored is indistinguishable from noise
unless the automaton measures it. Every term in the model's emergent
vocabulary therefore comes paired with an observable---shell completeness,
$u^2+v^2-R^4$, the fitted slope of $\Delta N$ versus $N$. This pairing is,
we believe, the methodological core of any serious ``it from bit''
programme.

Third, \emph{the honest objection}. Why this rule? A generic lattice
automaton does not look like a sphere, a polarization pair, or an attractor;
the design choices here are precisely the ones that produce those features.
This objection is not an argument against ``it from bit''---it \emph{is} the
research programme the slogan names. If physics is computation, the physical
laws are a program, and programs have authors; the question ``which program,
and why'' is exactly the question Wheeler bequeathed to us.

\section{Conclusion}
\label{sec:conclusion}

We have presented a complete, running automaton---a toy universe---in which a
spherical light clock, a momentum axis, a transverse polarization pair,
island-affine charge, and a matter--antimatter census all derive from a
single bit-local rule, without any of them being written into the initial
conditions as continuous fields. The construction demonstrates, at least in
miniature, the two theses that make ``it from bit'' more than a slogan: that
the standard-model vocabulary \emph{can} be instantiated by a purely
combinatorial system, and that emergence must be measured to be believed. The
instrumentation is the arbiter of the second thesis, and the choice of rule
is the open frontier of the first.

\begin{thebibliography}{99}

\bibitem{wheeler1990}
J.~A.~Wheeler, ``Information, physics, quantum: The search for links,''
in \emph{Complexity, Entropy, and the Physics of Information}, SFI Studies in
the Sciences of Complexity, vol.~VIII, W.~H.~Zurek, Ed. Redwood City, CA:
Addison-Wesley, 1990, pp.~3--28.

\bibitem{wheeler1989}
J.~A.~Wheeler, ``It from bit,'' in \emph{Proceedings of the 3rd International
Symposium on Foundations of Quantum Mechanics}, Tokyo, Japan, 1989,
pp.~354--368.

\bibitem{zuse1967}
K.~Zuse, ``Rechnender Raum,'' \emph{Elektronische Datenverarbeitung}, vol.~8,
pp.~336--344, 1967. English translation: ``Calculating space,'' in
\emph{A Computable Universe} \cite{zenil2012}.

\bibitem{fredkin1990}
E.~Fredkin, ``Digital mechanics,'' \emph{Physica D}, vol.~45, pp.~254--270,
1990.

\bibitem{wolfram2002}
S.~Wolfram, \emph{A New Kind of Science}. Champaign, IL: Wolfram Media, 2002.

\bibitem{toothooft2016}
G.~'t~Hooft, \emph{The Cellular Automaton Interpretation of Quantum
Mechanics}, Fundamental Theories of Physics, vol.~185. Cham: Springer, 2016.

\bibitem{zenil2012}
H.~Zenil, Ed., \emph{A Computable Universe: Understanding and Exploring
Nature as Computation}. Singapore: World Scientific, 2012.

\bibitem{vonneumann1966}
J.~von~Neumann, \emph{Theory of Self-Reproducing Automata},
A.~W.~Burks, Ed. Urbana, IL: University of Illinois Press, 1966.

\bibitem{lloyd2002}
S.~Lloyd, ``Computational capacity of the universe,'' \emph{Physical Review
Letters}, vol.~88, no.~23, art.~237901, 2002.

\bibitem{tegmark2008}
M.~Tegmark, ``The mathematical universe,'' \emph{Foundations of Physics},
vol.~38, no.~2, pp.~101--150, 2008.

\bibitem{toffoli1987}
T.~Toffoli and N.~Margolus, \emph{Cellular Automata Machines}. Cambridge,
MA: MIT Press, 1987.

\end{thebibliography}

\end{document}
